\subsection{フォーマット文字列攻撃}

\TT{puts(string)}や\TT{printf("\%s", string)}の代わりに\TT{printf(string)}を書くのはよくある間違いです。
攻撃者が自分のテキストを\TT{string}に入れることができれば、プロセスをクラッシュさせたり、ローカルスタック内の変数に洞察を得たりすることができます。

これを見てください：

\lstinputlisting[style=customc]{\CURPATH/f.c}

\printf が単一のフォーマット文字列以外に追加の引数を持たないことに注意してください。

攻撃者が最後の\printf の最初の引数に\TT{\%s}文字列を入れたと想像してみましょう。
このサンプルをx86 UbuntuでGCC 5.4.0を使ってコンパイルし、実行すると結果の実行ファイルは\q{world}文字列を出力します！

最適化をオンにすると、\printf はガベージを出力しますが―おそらく strcpy() 呼び出しが最適化されたか、ローカル変数も同様に最適化されたのでしょう。
また、x64コード、異なるコンパイラ、\ac{OS}などでは結果が異なります。

さて、攻撃者が\printf 呼び出しに次の文字列を渡すことができるとしましょう：\TT{\%x \%x \%x \%x \%x}。
私の場合、出力は：\q{80485c6 b7751b48 1 0 80485c0}（これらはローカルスタックからの値です）。
1と0の値、いくつかのポインタがあることがわかります（最初のものはおそらく\q{world}文字列へのポインタです）。
そのため、攻撃者が\TT{\%s \%s \%s \%s \%s}文字列を渡すと、プロセスはクラッシュします。なぜなら、\printf は1や0を文字列へのポインタとして扱い、そこから文字を読もうとして失敗するからです。

さらに悪いことに、コード内に\TT{sprintf (buf, string)}があり、\TT{buf}が1024バイト程度のローカルスタック内のバッファである場合、攻撃者は\TT{buf}がオーバーフローするように\TT{string}を細工でき、場合によってはコード実行につながる可能性があります。

多くの人気で有名なソフトウェアが（または今でも）脆弱でした：

\myindex{Quake}
\myindex{John Carmack}
\begin{framed}
\begin{quotation}
QuakeWorldが開始され、約4000ユーザーに達したが、その後マスターサーバーが爆発した。

DisrupterとChortsはより堅牢なコードに取り組んでいる。

もし誰かが故意にやったなら、教えてもらえませんか…（名前として\%sを試したすべての人ではありませんでした）
\end{quotation}
\end{framed}
( John Carmackの.planファイル、1996年12月17日\footnote{\url{https://github.com/ESWAT/john-carmack-plan-archive/blob/33ae52fdba46aa0d1abfed6fc7598233748541c0/by_day/johnc_plan_19961217.txt}} )

現在では、ほぼすべてのまともなコンパイラがこれについて警告します。

もう一つの問題は、あまり知られていない\TT{\%n} \printf 引数です：\printf がフォーマット文字列でそれに達すると、対応する引数にこれまでに出力された文字数を書き込みます：
\url{http://stackoverflow.com/questions/3401156/what-is-the-use-of-the-n-format-specifier-in-c}。
したがって、攻撃者はフォーマット文字列に多くの\TT{\%n}コマンドを渡すことでローカル変数を破壊できます。